\section{INFINITE SUMS AND PRODUCTS OF REAL NUMBERS}

Since every point of \(\mathbf{R}\) has a countable fundamental system of neighbourhoods (§ 1, no. 4, Corollary to Proposition 3), it follows that a family \((x_{t})\) of finite real numbers is summable in \(\mathbf{R}\) only if the set of indices \(t\) such that \(x_{t} \neq 0\) is countable (Chapter III, § 5, no. 2, Corollary to Proposition 1). The study of summable families in \(\mathbf{R}\) thus reduces essentially to the study of summable sequences. However, we shall later have to consider uncountable families \((x_{t})\) of finite real numbers, whose terms are functions of a parameter \(t\) ; it can happen that this family is summable for all t, but that the (countable) set of indices \(\iota\) such that \(x_{t} \neq 0\) depends on t. For this reason we shall not impose any hypothesis on the power of the index set in what follows.

